%% Q1
Learning is most effective when the \textbf{content (``what'') matches the method and style (``how'')} by which students learn, i.e.\ when teaching suits the learner's needs, level and learning style.
\why{Accurate presentation, time limits or completing the syllabus are teacher-centred concerns and do not by themselves guarantee learning.}

%% Q2
A profession is marked by specialised knowledge and skills, formal training institutions, a code of ethics and a service motive. A \textbf{market-oriented (profit-seeking) effort} is a feature of business, not a justification of teaching as a profession.

%% Q3
Beginning teachers typically struggle with content mastery, use of ICT and classroom management (Veenman's study of novice teachers). \textbf{Following institutional policy} is routine and is not a significant challenge.

%% Q4
Learning is said to be ``wasted'' when subjects are taught in an abstract, bookish way that is \textbf{not related to real life}. Students then cannot see the relevance of what they learn and do not retain or use it.

%% Q5
Using the local language makes the message understandable both when the teacher \textbf{encodes} it (puts ideas in familiar words) and when students \textbf{decode} it (understand it easily). It improves \textbf{encoding and decoding} together.

%% Q6
A teacher-organised \textbf{discussion} involves all students in speaking, listening, questioning and reasoning, so it ensures the greatest participation.
\why{Lecture and class notes are one-way; a field trip is occasional and not a routine classroom activity.}

%% Q7
Research begins with a \textbf{well-formulated problem}; hypotheses, literature review and design all follow from it. ``A problem well stated is half solved'' (Dewey).

%% Q8
Students across Gujarat form a large, geographically spread population. \textbf{Multistage random sampling} (e.g.\ districts $\rightarrow$ universities $\rightarrow$ colleges $\rightarrow$ students) gives a representative sample efficiently.

%% Q9
\textbf{Fundamental (basic) research} generates new knowledge, principles and theories, forming the foundation for later applied work.

%% Q10
The scientific method rests on \textbf{empirical evidence}, \textbf{experimentation} and \textbf{objectivity}. Mere \textbf{opinions and responses of people} are subjective and are not the basis of the scientific method.

%% Q11
The literature review helps identify research gaps, frame research questions and formulate hypotheses. The \textbf{implications of research} come from the researcher's \emph{own findings} at the end of the study, not from the review.

%% Q12
Continuous data are shown with a histogram, frequency polygon or ogive (cumulative frequency curve). A \textbf{bar chart}, with gaps between bars, suits \emph{discrete or categorical} data and is unsuitable for continuous data.

%% Q13
From smallest to largest: \textbf{KB $<$ MB $<$ GB $<$ TB}.

%% Q14
Routers, switches and hubs are \textbf{networking devices}. The \textbf{CPU} is the processing unit of a computer, not a network device.

%% Q15
Campus community radio (licensed since 2004) aims at educational and developmental broadcasting, i.e.\ \textbf{exposing students and the community to further information}.

%% Q16
Windows, Linux and macOS are \textbf{operating systems}. \textbf{Google} is a company and search engine, so it is the odd one out. (Google's operating systems are Android and ChromeOS.)

%% Q17
Genuine (two-way) communication is dynamic. Each message depends on the \textbf{spontaneous verbal and non-verbal feedback} of the other person, so it cannot be fully planned in advance.

%% Q18
\textbf{HTTP} is \textbf{HyperText Transfer Protocol}, the protocol of the World Wide Web.

%% Q19
Communication is successful only when there is shared meaning: \textbf{both sender and receiver understand} the information being transferred.

%% Q20
Video conferencing carries \textbf{sound and pictures} in \textbf{both directions} in real time, so it is audio-visual two-way communication.

%% Q21
VoIP offers low-cost calls, portability (one number anywhere) and extra features. \textbf{``Increased geographical anonymity''} is not normally regarded as an advantage; it complicates emergency services and legal tracing.

%% Q22
Communication within oneself (thinking, self-talk, reflection) is \textbf{intrapersonal communication}.

%% Q23
\textbf{Bluetooth} is a short-range \textbf{wireless} technology (2.4 GHz band, about 10 m) for connecting devices such as phones, headsets, keyboards and speakers.

%% Q24
Communication is effective when the message \textbf{reaches the receiver as intended by the sender}, i.e.\ the intended meaning is understood.

%% Q25
Each term is (previous $\times2$) $+1$: $5\rightarrow11\rightarrow23\rightarrow\mathbf{47}\rightarrow95\rightarrow191$.

%% Q26
ARUTF $\rightarrow$ DSXUI: the shifts are $+3, +1, +3, +1, +3$.
PQRSF $\rightarrow$ P$+3=$S, Q$+1=$R, R$+3=$U, S$+1=$T, F$+3=$I, giving \textbf{SRUTI}.

%% Q27
SREESANTH $\rightarrow$ RSDFRBMUG: the shifts alternate $-1, +1, -1, +1, \ldots$
FIXING $\rightarrow$ F$-1=$E, I$+1=$J, X$-1=$W, I$+1=$J, N$-1=$M, G$+1=$H, giving \textbf{EJWJMH}.

%% Q28
The hands are opposite when they are $180^\circ$ apart. The angle is $|30H-5.5M|$, so $|120-5.5M|=180$ gives $5.5M=300$ and
\[M=\frac{300}{5.5}=\frac{600}{11}=54\tfrac{6}{11}\text{ minutes past 4}\]

%% Q29
PEN and HEN differ only in the first letter: P $=16$, H $=8$, and $25-17=8$ matches $16-8$. So the value is (position of first letter) $+9$: M $=13$, giving $13+9=\mathbf{22}$.

%% Q30
Bindu is 12th from each end: $11$ girls on each side $+1=\mathbf{23}$.

%% Q31 ctx
\textbf{Read the following text and answer the questions from 31 to 36.}\par\smallskip
Investigators have noticed some similarities between the fantastic events in fairy tales and those in adult dreams --- the fulfilment of wishes, winning over all competitors and the destruction of enemies. There are, of course, very significant differences also. For example, in dreams more often than not the wish fulfilment is disguised, while in fairy tales much of it is openly expressed. To a considerable degree, dreams are the result of inner pressures which have found no relief, and of problems which beset a person to which he knows no solution and to which the dream finds none. The fairy tale does the opposite; it projects the relief of all pressures and not only offers ways to solve problems, but promises a happy solution. We cannot control what goes on in our dreams. Although our inner censorship influences what we may dream, such control occurs at an unconscious level. The fairy tale, on the other hand, is the result of conscious and unconscious content having been shaped by the conscious mind, not of one particular person but of many in regard to what they view as universal human problems. Only if a fairy tale met the conscious and unconscious needs of many people, would it be told and retold with great interest.

%% Q31
The author discusses wishes, inner pressures, the conscious and the unconscious, i.e.\ the \textbf{psychological relevance} of fairy tales (the passage echoes Bruno Bettelheim's \emph{The Uses of Enchantment}).

%% Q32
The similarities listed are ``the fulfilment of wishes, winning over all competitors and the \textbf{destruction of enemies}''.

%% Q33
Fairy tales differ from dreams in being shaped by \textbf{many people} (communal creation, I) and in promising a \textbf{happy solution} (II). Enduring appeal (III) is a consequence, not a stated point of difference. The answer is (I) and (II) only.

%% Q34
The passage says only a tale that meets the needs of many people is retold with great interest, i.e.\ \emph{not all} tales. So ``all tales are retold with great interest'' is \textbf{not true}.

%% Q35
The passage sets out similarities and then differences between dreams and fairy tales: \textbf{comparison and contrast}.

%% Q36
The official key gives \textbf{Psychology of dreams}: the passage analyses wish fulfilment, inner pressures, censorship and the conscious and unconscious mind, which are psychological aspects of dreams (with fairy tales used for contrast).

%% Q36 note
Many readers would find ``Dreams and fairy tales'' (B) a closer title, since the passage compares both throughout. The official key is (A).

%% Q37
Passive (second-hand) smoke \textbf{harms the lungs and health of non-smokers}. This health hazard is the basis of the ban under the Cigarettes and Other Tobacco Products Act (COTPA), 2003, and the 2008 rules on smoking in public places.

%% Q38
Waste management follows ``reduce, reuse, \textbf{recycle}''. \textbf{Recycling biodegradable waste} (composting, biogas) is a direct example.

%% Q39
Graffiti, loud audio and noisy motorbikes harm the environment (visual and noise pollution). \textbf{Going on strike} is a social or political action, not an environmental one.

%% Q40
Non-conventional sources (solar, wind, biogas) replace fossil fuels and so reduce the \textbf{consumption of crude oil}.

%% Q41
The \textbf{Narmada Bachao Andolan} (Medha Patkar, Baba Amte) arose from the \textbf{displacement and rehabilitation} of people affected by the Sardar Sarovar and other Narmada dams, along with environmental concerns.

%% Q42
Seasons are caused by the \textbf{tilt of the Earth's axis} (about $23.5^\circ$) as the Earth revolves around the Sun. This changes the angle of sunlight and the length of the day.
\why{Rotation causes day and night.}

%% Q43
(i) ``Union elections should not be held'' is an \textbf{opinion}. (ii) ``Student leaders are never able to perform well'' is a \textbf{prejudice} (an unfair generalisation). (iii) ``Do study regularly'' is \textbf{advice}. (iv) ``Degree cannot be awarded without credits'' is a \textbf{fact}.
Opinion, Prejudice, Fact, Advice $=$ (i), (ii), (iv), (iii), which is option (B).

%% Q44 opts
\tikz[baseline=-0.5ex,scale=0.3]{\draw (0,0) circle(1.2); \draw (0,0) circle(0.6); \draw (2.2,0) circle(0.7);}
\tikz[baseline=-0.5ex,scale=0.3]{\draw (0,0) circle(1); \draw (2.2,0) circle(0.75); \draw (3.8,0) circle(0.5);}
\tikz[baseline=-0.5ex,scale=0.3]{\draw (0,0) circle(1.1); \draw (-1.2,0.15) circle(0.6); \draw (1.2,0.15) circle(0.6);}
\tikz[baseline=-0.5ex,scale=0.3]{\draw (0,0) circle(1.3); \draw (-0.55,0) circle(0.5); \draw (0.55,0) circle(0.5);}

%% Q44
The official key gives diagram (C): a large circle (\textbf{travellers}) partly overlapping two separate circles (\textbf{train} and \textbf{bus}). Some travellers use the train and some the bus, while the train and bus sets also contain elements that are not travellers (crew, vehicles). No one is both in a train and in a bus at the same time.
\tip{Compare GSET Nov 2022 Q.~45, where the key showed train and bus inside the travellers circle. Read each figure carefully.}

%% Q45 note
Strictly, the pair that cannot both be true but can both be false (\emph{contraries}) is (i) and (iii), which is not offered. The official key gives (D).

%% Q45
Contrary propositions (``All S are P'' and ``No S is P'') cannot both be true but can both be false. Here those are (i) and (iii). (ii) and (iii) are \emph{contradictories} (exactly one is true). The key's answer (D) pairs (i) with (iv), and these two statements are about different predicates. Mark (D) as per the key.

%% Q46
``Bouncing up after hitting the bottom'' assumes that a person can hit the bottom, i.e.\ fail completely (i), and can rise again after failure (ii). \textbf{Both} assumptions are implicit.

%% Q47
Runs scored say nothing about how many players were spinners (i) or who opened the batting (ii). \textbf{Neither} conclusion follows.

%% Q48
P and Q always have the same truth value (P true $\Leftrightarrow$ Q true), so they are \textbf{logically equivalent}.

%% Q49
Readers who choose to respond select themselves, so this is a \textbf{self-selected (voluntary response)} sample, which is prone to bias.

%% Q50 ctx
A survey asked people how often they exceed speed limits. The data is then categorized into the following table of counts showing the relationship between age group and response. Read the table and answer the Question Nos.~50 to 54.
\begin{center}
\begin{tabular}{lccc}\toprule
& \multicolumn{2}{c}{\textbf{Exceed limit if possible}} & \\\cmidrule(lr){2-3}
\textbf{Age} & Always & Not always & \textbf{Total}\\\midrule
Under 30 & 100 & 100 & 200\\ Over 30 & 40 & 160 & 200\\\midrule Total & 140 & 260 & 400\\\bottomrule
\end{tabular}
\end{center}

%% Q50
Odds $=$ always : not always $=100:100=\mathbf{1\ to\ 1}$.
\tip{Odds compare \emph{for} with \emph{against}; risk (probability) compares \emph{for} with the \emph{total}.}

%% Q51
The positive attitude is ``not always'' exceeding. Over 30: $160/200=80\%$; under 30: $100/200=50\%$. Difference $=\mathbf{30\%}$.

%% Q52
Overall risk $=\dfrac{140}{400}=\mathbf{0.35}$.

%% Q53
Risk under 30 $=100/200=0.5$; risk over 30 $=40/200=0.2$.
\[\text{Relative risk}=\frac{0.5}{0.2}=\mathbf{2.5}\]

%% Q54
Risk over 30 $=\dfrac{40}{200}=\mathbf{0.20}$.

%% Q55
In an \textbf{open learning} system entry barriers are removed. Per the key, the student \textbf{can join any programme at any time}, with flexibility in admission, pace, place and choice of courses.

%% Q56
The \textbf{impact factor} (Eugene Garfield; Clarivate's Journal Citation Reports) is a measure for \textbf{journals}: the average number of citations received in a year by articles the journal published in the previous two years.

%% Q57
Maulana Abul Kalam Azad's birthday, \textbf{11 November}, is celebrated as \textbf{National Education Day} (since 2008).

%% Q58
In a credit system, \textbf{one credit $=$ one hour of teaching (lecture or tutorial) per week} over a semester, or two hours of practical per week.

%% Q59
A private university in India is established by an \textbf{Act of the State Legislature} (in Gujarat, under the Gujarat Private Universities Act, 2009) and must conform to the UGC (Private Universities) Regulations, 2003.

%% Q60
NAAC accreditation certifies the \textbf{level of quality} of an institution, expressed as a grade (A++ to C) based on the CGPA.
